\section{Local fundamental groups and exceptional lattices}\label{sec:local}
In this section all punctured local schemes are formed after strict henselization.

For a closed point $z$ of a normal surface $Z$, write $U_z^\circ = \Spec\O^{\rm sh}_{Z,z}\setminus\{z\}$ and let $\pi_1^{\rm et}(U_z^\circ)$ denote its \'etale fundamental group.
We recall the following standard comparison lemma.

\begin{lemma}\label{lem:transport}
Let $k$ be a separably closed field, let $X$ be a normal surface over $k$, and let $\nu\from S\to (X_{\bar k})_{\rm red}$ be the normalization, with induced morphism $\phi\from S\to X$.
Let $s\in S$ be a closed point with image $x\in X$.
Then $\phi$ is a universal homeomorphism and induces a universal homeomorphism $U_s^\circ\to U_x^\circ$, hence an isomorphism $\pi_1^{\rm et}(U_s^\circ)\cong\pi_1^{\rm et}(U_x^\circ)$.
If $X$ is regular at $x$, then $\pi_1^{\rm et}(U_s^\circ) = 1$.
\end{lemma}

\begin{proof}
Since $\bar k/k$ is purely inseparable, $\phi$ is a universal homeomorphism by \cite[Lemma~2.2(2)]{Tan18BC}.
The scheme $S\times_X\Spec\O^{\rm sh}_{X,x}$ is strictly henselian and ind-\'etale over $S$, hence isomorphic to $\Spec\O^{\rm sh}_{S,s}$.
Removing the closed points gives a universal homeomorphism $U_s^\circ\to U_x^\circ$, so topological invariance gives the asserted isomorphism \cite[Tag~\stackstag{04DY}]{Stacks}.
The regular-point case follows from \cite[Corollary~2.4]{Sch08}.
\end{proof}

We next construct a regular double cover of the singularity produced by the contraction in Lemma~\ref{lem:contraction}.

\begin{lemma}\label{lem:A1-cover}
Let $k$ be a separably closed field of characteristic $p>3$, let $X$ be a regular projective surface over $k$ with $K_X\equiv 0$, and let $f\from X\to X'$ be the contraction of a negative curve $C$ as in Lemma~\ref{lem:contraction}.
Let $S$ be the normalization of $(X'_{\bar k})_{\rm red}$ and let $s\in S$ lie over $P = f(C)$.
Then $\pi_1^{\rm et}(U_P^\circ)\cong\pi_1^{\rm et}(U_s^\circ)\cong\ZZ/2$.
\end{lemma}

\begin{proof}
Put $k_C = H^0(C,\O_C)$ and $A = \O^{\rm sh}_{X',P}$, a strictly henselian local ring with residue field $k_C$.
Define $T$ by the Cartesian square
\[
  \begin{tikzcd}
    T \arrow[r] \arrow[d,"g"'] & X \arrow[d,"f"]\\
    \Spec A \arrow[r] & X'.
  \end{tikzcd}
\]
The projection $T\to X$ is ind-\'etale, so $T$ is regular \cite[Tags~\stackstag{07EP} and~\stackstag{0H7S}]{Stacks}, with exceptional curve $C\cong\PP^1_{k_C}$.

We first construct a double cover of $T$.
Put $T_0 = T\times_A k_C$.
By proper base change \cite[Exp.~XII, Corollaire~5.5(iii)]{SGA4}, $H^2_{\rm et}(T,\mu_2)\cong H^2_{\rm et}(T_0,\mu_2)$.
Since $C = (T_0)_{\rm red}$, restriction induces $H^2_{\rm et}(T_0,\mu_2)\cong H^2_{\rm et}(C,\mu_2)$.
After the purely inseparable extension $\overline{k_C}/k_C$, the standard cohomology calculation for $C_{\overline{k_C}}\cong\PP^1_{\overline{k_C}}$ gives $H^2_{\rm et}(C,\mu_2)\cong\ZZ/2$ \cite[Tag~\stackstag{03RQ}]{Stacks}.
The Kummer sequence induces an exact sequence \cite[Tag~\stackstag{03PL}]{Stacks}
\begin{equation}\label{eq:kummer-exact-sequence}
  \Pic(T)\xrightarrow{(-)^{\otimes 2}}\Pic(T)
  \xrightarrow{\delta_T}H^2_{\rm et}(T,\mu_2).
\end{equation}
For a line bundle $\mathcal L$ on $T$, functoriality of the Kummer sequence shows that the image of $\delta_T(\mathcal L)$ under $H^2_{\rm et}(T,\mu_2)\xrightarrow{\sim}H^2_{\rm et}(C,\mu_2)\cong\ZZ/2$ is $\deg_{k_C}(\mathcal L|_C)\bmod 2$.
Since $\deg_{k_C}(\O_T(C)|_C) = -2$, exactness of \eqref{eq:kummer-exact-sequence} gives a line bundle $M$ on $T$ with $M^{\otimes 2}\cong\O_T(C)$.
Let
\[
  \varpi\from V = \operatorname{\mathbf{Spec}}_T(\O_T\oplus M^{-1})\to T
\]
be the double cover obtained locally by adjoining a square root of a local equation $t = 0$ for $C$.
At a closed point $x\in C$, regularity of $T$ and $C$ allows us to choose regular parameters $t,u$ in $\O_{T,x}$.
The local ring of $V$ above $x$ has dimension two and maximal ideal $(z,u)$, since $t = z^2$, so it is regular.
Over $T\setminus C$, the derivative $2z$ is invertible, so the cover is \'etale and $V$ is regular there.
The cover is integral because $t$ has order one along $C$.
The ramification curve $D = \varpi^{-1}(C)_{\rm red}$ satisfies $D\cong C$ and $2D = \varpi^*C$.
Thus $\O_V(D)|_D\cong\varpi^*M|_D$ has degree $-1$ over $k_C$, giving $\O_D(-jD)\cong\O_{\PP^1_{k_C}}(j)$ for $j\geq 0$.

We next show that $B = H^0(V,\O_V)$, the normalization of $A$ in the function field of $V$, is regular.
Since $B$ is finite over the henselian ring $A$ and is a domain, it is local and henselian.
The closed fibre of $V\to\Spec B$ is supported on $D$.
Since $-C$ is ample over $\Spec A$ and $2D = \varpi^*C$, the divisor $-D$ is ample over $\Spec B$.
Relative Serre vanishing and descending induction in the exact sequences
\[
  0\to\O_V(-(j+1)D)\to\O_V(-jD)\to\O_D(-jD)\to 0
\]
give $H^1(V,\O_V(-jD)) = 0$ for all $j\geq 0$, since $H^1(D,\O_D(-jD)) = 0$.
By taking cohomology of the above sequence with $j = 0$, the restriction $B\to H^0(D,\O_D) = k_C$ is surjective, so the residue field of $B$ is $k_C$.
By the same argument as in the proof of Lemma~\ref{lem:contraction}\,(5), we have
\[
  \operatorname{gr}_{\mathfrak m_B}\widehat B
  \cong\bigoplus_{j\geq 0}H^0(\PP^1_{k_C},\O(j))
  \cong k_C[u,v].
\]
Hence $B$ is a regular local ring of dimension two and is strictly henselian.

Finally, we compute the local fundamental groups.
Over $U_P^\circ$, the finite morphism $\Spec B\to\Spec A$ is the \'etale double cover $V\setminus D\to T\setminus C$.
Put $U_B^\circ = \Spec B\setminus\{\mathfrak m_B\}$.
Then $U_B^\circ\to U_P^\circ$ is a connected \'etale double cover.
By purity \cite[Exp.~X, Corollaire~3.3]{SGA1}, $\pi_1^{\rm et}(U_B^\circ)\cong\pi_1^{\rm et}(\Spec B) = 1$, since $B$ is strictly henselian.
It follows that $\pi_1^{\rm et}(U_P^\circ)\cong\ZZ/2$.
Lemma~\ref{lem:transport} gives $\pi_1^{\rm et}(U_s^\circ)\cong\pi_1^{\rm et}(U_P^\circ)$.
\end{proof}

The next lemma relates the local fundamental group to the exceptional curves and their intersection matrix, without assuming that the singularity is rational.
It is an \'etale analogue of Mumford's computation of the first homology of the link of a complex normal surface singularity \cite[Part~I, pp.~10--11]{Mum61}.
In particular, it restricts the possible intersection matrices in the regular-point case of Lemma~\ref{lem:transport} and in the setting of Lemma~\ref{lem:A1-cover}.

\begin{lemma}\label{lem:kummer}
Let $S$ be a normal surface over an algebraically closed field of characteristic $p$, let $s\in S$ be a closed point, let $A\supseteq\O^{\rm sh}_{S,s}$ be a normal local domain finite over $\O^{\rm sh}_{S,s}$, and let $g\from T\to\Spec A$ be a resolution of singularities.
Let $E = E_1\cup\dots\cup E_r$ be the reduced exceptional curve of $g$, let $M = (E_i\cdot E_j)_{\kappa(s)}$, and let $U_s^\circ = \Spec A\setminus\{\mathfrak m_A\}$.
The following hold for every positive integer $n$ prime to $p$.
\begin{enumerate}[label=\textup{(\arabic*)}]
\item There is an exact sequence of \'etale cohomology groups
\begin{equation}\label{eq:kummer-lattice-sequence}
  0\to H^1_{\rm et}(E,\mu_n)\to H^1_{\rm et}(U_s^\circ,\mu_n)\to\ker\bigl(M\otimes\ZZ/n\from(\ZZ/n)^r\to(\ZZ/n)^r\bigr)\to 0.
\end{equation}
An invertible sheaf $\mathcal L$ on $T$ has an $n$-th root in $\Pic(T)$ if and only if $n\mid\deg(\mathcal L|_{E_i})$ for every $i$.
\item If $\pi_1^{\rm et}(U_s^\circ) = 1$, then $H^1_{\rm et}(E,\mu_n) = 0$, every $E_i$ has geometric genus zero, and $\lvert\det M\rvert$ is a power of $p$.
\item If $p\neq 2$, $\pi_1^{\rm et}(U_s^\circ)\cong\ZZ/2$ and $H^1_{\rm et}(E,\mu_4) = 0$, then $\lvert\det M\rvert = 2p^e$ for some $e\geq 0$.
The vanishing assumption holds if $E$ is a tree of smooth rational curves meeting transversally.
\end{enumerate}
\end{lemma}

\begin{proof}
For \textup{(1)}, the ring $A$ is strictly henselian with residue field $\kappa(s)$.
Hensel's lemma shows that every unit of $A$ is an $n$-th power, so $A^*$ is $n$-divisible.
Normality gives $H^0(T,\O_T) = H^0(U_s^\circ,\O_{U_s^\circ}) = A$.
As in the proof of Lemma~\ref{lem:A1-cover}, we have $H^i_{\rm et}(T,\mu_n)\cong H^i_{\rm et}(E,\mu_n)$ for $i = 1,2$.
Let $\pi\from\widetilde E\to E$ be the normalization, with components $\widetilde E_i$ over the $E_i$.
Since $\pi$ is finite and the cokernel of $\mu_n\to\pi_*\mu_n$ is supported at finitely many points, the normalization sequence gives a surjection $H^1_{\rm et}(E,\mu_n)\to H^1_{\rm et}(\widetilde E,\mu_n)$ and an isomorphism $H^2_{\rm et}(E,\mu_n)\cong H^2_{\rm et}(\widetilde E,\mu_n)\cong(\ZZ/n)^r$ \cite[Tag~\stackstag{03RQ}]{Stacks}.
For $Y = T$ or $Y = U_s^\circ$, we have $H^0(Y,\O_Y)^* = A^*$, so the Kummer sequence gives
\[
  0\to A^*/(A^*)^n\to H^1_{\rm et}(Y,\mu_n)\to\Pic(Y)[n]\to 0.
\]
Since $A^*$ is $n$-divisible, the quotient $A^*/(A^*)^n$ vanishes, and we obtain
\begin{equation}\label{eq:kummer-identifications}
  \begin{aligned}
    \Pic(T)[n] &\cong H^1_{\rm et}(T,\mu_n)\cong H^1_{\rm et}(E,\mu_n),\\
    \Pic(U_s^\circ)[n] &\cong H^1_{\rm et}(U_s^\circ,\mu_n).
  \end{aligned}
\end{equation}
The Kummer boundary map $\Pic(T)\to H^2_{\rm et}(T,\mu_n)$ has kernel $n\Pic(T)$, so it induces an injection $\Pic(T)/n\embed H^2_{\rm et}(T,\mu_n)\cong(\ZZ/n)^r$.
Since the Kummer boundary map commutes with pullback, the image of $[\mathcal L]$ has $i$-th coordinate $\deg(\mathcal L|_{E_i})\bmod n$.
Thus $\mathcal L$ has an $n$-th root if and only if all these coordinates vanish.
This proves the assertion about $n$-th roots of line bundles.

Every line bundle on $U_s^\circ$ is represented by a divisor whose closure is Cartier on $T$, since $T$ is regular.
Thus the restriction map $\operatorname{res}\from\Pic(T)\to\Pic(U_s^\circ)$ is surjective.
For $[\mathcal L]\in\ker(\operatorname{res})$, a trivialization of $\mathcal L|_{U_s^\circ}$ defines a rational section of $\mathcal L$ whose divisor is supported on $E$.
Thus $\ker(\operatorname{res})$ consists of the classes $\O_T(\sum_i a_iE_i)$ with $a_1,\ldots,a_r\in\ZZ$.
If $\O_T(\sum_i a_iE_i)$ is trivial, then the negative definiteness of $M$ forces all $a_i = 0$.
We therefore have an exact sequence
\[
  0\to\ZZ^r\to\Pic(T)\xrightarrow{\operatorname{res}}\Pic(U_s^\circ)\to 0,
\]
whose first map sends $(a_i)_i$ to $\O_T(\sum_i a_iE_i)$.
The long exact sequence for $\operatorname{Tor}^{\ZZ}_i(-,\ZZ/n)$ gives $0\to\Pic(T)[n]\to\Pic(U_s^\circ)[n]\to\ker\bigl((\ZZ/n)^r\to\Pic(T)/n\bigr)\to 0$.
The map $(\ZZ/n)^r\to\Pic(T)/n$ sends the reduction of $(a_i)_i\in\ZZ^r$ to $[\O_T(\sum_i a_iE_i)]$.
For each $i$, the degree of this line bundle on $E_i$ is
\[
  \deg\bigl(\O_T(\textstyle\sum_j a_jE_j)|_{E_i}\bigr)
  = \sum_j(E_i\cdot E_j)a_j.
\]
The composite $(\ZZ/n)^r\to\Pic(T)/n\embed(\ZZ/n)^r$ is therefore $M\otimes\ZZ/n$.
Composing with the injection constructed above preserves the kernel.
Together with the Kummer identifications in \eqref{eq:kummer-identifications}, this proves \eqref{eq:kummer-lattice-sequence}.

To prove \textup{(2)} and \textup{(3)}, first note that choosing a primitive $n$-th root of unity in the ground field identifies $\mu_n$ with the constant sheaf $\ZZ/n$ on $U_s^\circ$.
As $U_s^\circ$ is connected, we have $H^1_{\rm et}(U_s^\circ,\mu_n)\cong\operatorname{Hom}(\pi_1^{\rm et}(U_s^\circ),\ZZ/n)$.

For \textup{(2)}, the exact sequence \eqref{eq:kummer-lattice-sequence} gives $H^1_{\rm et}(E,\mu_n) = 0$ and $\ker(M\otimes\ZZ/n) = 0$ for every such $n$.
For a prime $\ell\neq p$, the normalization surjection then gives $H^1_{\rm et}(\widetilde E_i,\mu_\ell) = 0$, so each $E_i$ has geometric genus zero.
Moreover, $\ker(M\otimes\ZZ/\ell) = 0$ implies $\ell\nmid\det M$, proving that $\lvert\det M\rvert$ is a power of $p$.

For \textup{(3)}, if $\ell\neq p$ is an odd prime, then $\ker(M\otimes\ZZ/\ell)$ is a quotient of $\operatorname{Hom}(\ZZ/2,\ZZ/\ell) = 0$, so $\ell\nmid\det M$.
As $H^1_{\rm et}(E,\mu_4) = 0$, the group $\ker(M\otimes\ZZ/4)\cong\operatorname{Hom}(\ZZ/2,\ZZ/4)$ has order two.
Writing $d_1,\dots,d_r$ for the invariant factors of $M$, we have $\ker(M\otimes\ZZ/4)\cong\bigoplus_i\ZZ/\gcd(d_i,4)$.
Thus exactly one $d_i$ is even and it is not divisible by four, so $\det M$ has $2$-adic valuation one, proving the determinant formula.
If $E$ is a tree of smooth rational curves meeting transversally, then $\Pic(E)\cong\ZZ^r$ and $H^0(E,\O_E)^* = \kappa(s)^*$.
The latter group is $4$-divisible, so the Kummer sequence gives $H^1_{\rm et}(E,\mu_4)\cong\Pic(E)[4] = 0$.
\end{proof}

\begin{corollary}\label{lem:pi1-lattice}
Let $k$, $W$ and $\phi\from S\to W$ be as in Lemma~\ref{lem:transport}, let $s\in S$ be a singular point with image $w\in W$, and let $g\from T\to\Spec\O^{\rm sh}_{S,s}$ be a resolution with reduced exceptional curve $E = E_1\cup\dots\cup E_r$ and intersection matrix $M = (E_i\cdot E_j)$.
\begin{enumerate}
\item If $W$ is regular at $w$, then $H^1_{\rm et}(E,\mu_n) = 0$ for every integer $n$ prime to $p$, every $E_i$ has geometric genus zero, and $\lvert\det M\rvert = p^e$ for some $e\geq 0$.
\item If $s$ is a rational singularity and $g$ is minimal, then $E$ is a tree of smooth rational curves meeting transversally.
  If moreover $W = X'$ and $w = P$ are as in Lemma~\ref{lem:A1-cover}, then $\lvert\det M\rvert = 2p^e$ for some $e\geq 0$.
\end{enumerate}
\end{corollary}

\begin{proof}
Apply Lemma~\ref{lem:kummer} with $A = \O^{\rm sh}_{S,s}$.
Part~\textup{(1)} follows from Lemmas~\ref{lem:transport} and~\ref{lem:kummer}(2).
In \textup{(2)}, the description of $E$ follows from rationality \cite{Lip69}, and the determinant assertion follows from Lemmas~\ref{lem:A1-cover} and~\ref{lem:kummer}(3).
\end{proof}


Let $s$ be a rational double point on a normal surface $S$ over an algebraically closed field $k$ of characteristic $p$, and put $A = \O^{\rm sh}_{S,s}$.
Let $T\to\Spec A$ be the minimal resolution, with exceptional curves $E_1,\dots,E_r$ and intersection matrix $M$.
Table~\ref{tab:rdp-invariants} records the equations and the absolute determinants $\lvert\det M\rvert$.
The determinants follow from the ADE intersection matrices and are valid in every characteristic.
If $p>5$, then $\widehat A\cong k[[x,y,z]]/(f)$ with $f$ as in the table \cite[Section~3]{Art77}.

\begin{table}[htbp]
\centering
\small
\setlength{\tabcolsep}{4pt}
\renewcommand{\arraystretch}{1.3}
\begin{tabular}{@{}lll@{}}
\hline
Type & Equation $f$ & $\lvert\det M\rvert$ \\
\hline
$\mathrm A_n$, $n\geq 1$ & $x^{n+1}+yz$ & $n+1$ \\
$\mathrm D_n$, $n\geq 4$ & $z^2+x^2y+y^{n-1}$ & $4$ \\
$\mathrm E_6$ & $z^2+x^3+y^4$ & $3$ \\
$\mathrm E_7$ & $z^2+x^3+xy^3$ & $2$ \\
$\mathrm E_8$ & $z^2+x^3+y^5$ & $1$ \\
\hline
\end{tabular}
\caption{Equations for $p>5$ and exceptional-lattice determinants of rational double points.}
\label{tab:rdp-invariants}
\end{table}

Combining Lemma~\ref{lem:transport} and Corollary~\ref{lem:pi1-lattice}(1) with the known local fundamental groups of rational double points gives the following restriction.
Hirokado \cite[Theorem~3.4]{Hir04} proved a special case; the general form below is due to Schr\"oer.

\begin{corollary}[{\cite[Theorem~6.1]{Sch08}}]\label{lem:geometric-rdp}
Let $X$ be a regular surface over a field $k$ of characteristic $p>5$.
Then every rational double point of $X_{\bar k}$ has type $\mathrm A_{p^e-1}$ for some integer $e\geq 1$.
\end{corollary}
